%----------------------------------------------------------------------------------------
%	DESCRIPTION OF THE PRODUCTS THAT ARE BEING EXPECTED FROM THE STAGE
%----------------------------------------------------------------------------------------
\section*{Prodotti attesi}
% Personalizzare definendo i prodotti attesi (facoltativo)
Lo studente dovrà produrre i seguenti prodotti:
\begin{enumerate}
    \item \textbf{Codice sorgente AI Agent} \newline
    Implementazione del sistema software completo, comprendente: l'automazione Python per l'integrazione Elastic-AppDynamics, il nucleo dell'AI Agent per l'analisi delle metriche e il rilevamento delle anomalie, e il workflow automatizzato di gestione degli incidenti.

    \item \textbf{Prompt} \newline
    Insieme dei prompt progettati per l'interazione con i modelli LLM, ottimizzati per l'analisi delle metriche di traffico, la classificazione degli incidenti e la generazione di output strutturati (diagnosi, suggerimenti, classificazione).

    \item \textbf{Documentazione tecnica} \newline
    Documentazione completa del sistema, comprendente la documentazione tecnica (architettura, flussi decisionali, esempi di utilizzo) e la documentazione funzionale (use case).
\end{enumerate}

Nel caso in cui lo studente abbia ancora tempo a sua disposizione, potrà approfondire il confronto tra i principali modelli LLM (OpenAI, Anthropic Claude, Microsoft Copilot) in termini di qualità delle analisi, costi e latenza (obiettivo facoltativo F01).
